\documentclass[11pt]{article}

\usepackage[T1]{fontenc}
\usepackage[utf8]{inputenc}
\usepackage{amsmath,amssymb,amsthm}
\usepackage[margin=1.15in]{geometry}
\usepackage{xcolor}
\usepackage[colorlinks=true,linkcolor=blue!50!black,citecolor=blue!50!black,urlcolor=blue!50!black]{hyperref}
\expandafter\def\expandafter\UrlBreaks\expandafter{\UrlBreaks\do\-\do\_\do\.}

\newtheorem{proposition}{Proposition}
\newtheorem{lemma}{Lemma}[section]
\newtheorem{observation}[lemma]{Observation}
\theoremstyle{definition}
\newtheorem{definition}[lemma]{Definition}
\theoremstyle{remark}
\newtheorem{remark}[lemma]{Remark}
\newtheorem{question}[lemma]{Question}

\newcommand{\N}{\mathbb{N}}
\newcommand{\Z}{\mathbb{Z}}
\newcommand{\FS}{\operatorname{FS}}
\newcommand{\loc}[1]{\mathcal{L}_{#1}}
\DeclareMathOperator{\SSum}{\Sigma}

\title{Sharp local carry lemmas and a quadratic radix barrier for
admissible-set block constructions\\[2pt]
{\large (Erd\H{o}s Problem 875)}}
\author{\textsc{erdosSweep} harness\thanks{Machine-generated mathematics,
adversarially cross-verified; see the provenance statement in
Section~\ref{sec:provenance}. Harness maintainer: Douglas~R.~Colkitt.}}
\date{July 2026}

\begin{document}
\maketitle

\begin{abstract}
An infinite set $A\subset\N$ is \emph{admissible} if any two finite subsets
with equal sums have equal cardinalities. The strongest known bound on the
gaps of an admissible sequence, $a_{n+1}-a_n\le n^{3+2\sqrt2}$ (GPT-5.5
Pro--Mazur, 2026), rests on a local ``carry lemma'' for the block
$C_m=\{1+i(m+1):0\le i<m\}$ modulo $Q_m=m^2+m+1$, proved there for carry
radius $|u|<m/4$ under the quotient window $|w|<m/2$. We prove the sharp
local statement: the carry lemma for $(C_m,Q_m)$ holds at radius exactly
$\lceil m/2\rceil$ with no hypothesis on the quotient, and fails at radius
$\lceil m/2\rceil+1$. (At this radius, quotients realized by $C_m$
automatically satisfy $2|w|<m$, so the substance is the sharp radius,
not the absent window.) This upgrades the closure step of the construction: a block
extension preserves the carry invariant assuming only the cardinality bound
$|A|\le\lceil k/2\rceil$ --- no containment $A\subset[1,M)$ and no mass
estimate. In the opposite direction we delimit the framework. First, every
$k$-element block satisfying the exact scalar carry property \emph{at all
quotients} has radix $Q\ge\lfloor k^2/4\rfloor+1$, and
$Q\ge\lfloor(k+1)^2/4\rfloor+1$ when $Q\notin C$; both bounds are attained.
The subset-sum counting input is classical (Nathanson 1995; Bhanja--Pandey
2022); we include a short self-contained proof. This barrier applies
\emph{only} to the all-quotient property: invariants imposed for a bounded
window of quotients are not ruled out and remain the natural attack
direction. Second, staggered lifts $c_i\mapsto c_i+Q_m\lambda_i$ of the
Mazur block preserve the exact invariant precisely when
$\lambda_i=\lambda(m-i)$, and unit carry radius then forces $\lambda=0$;
this rigidity persists under quotient windows comparable to the lift size.
The sharp carry lemma and the closure --- including the sharpness witness,
for every $m$ --- have been formalized in Lean~4 with mathlib; the barrier
and rigidity results are paper-only.
\emph{This note resolves none of the questions of Erd\H{o}s Problem 875}:
it contains no new construction, no improved gap exponent, and no result on
the ratio question. It sharpens the local machinery of the known
construction and shows that further progress within the scalar-carry
framework must come from bounded-window invariants, non-scalar invariants,
or new residue systems.
\end{abstract}

\section{Introduction}\label{sec:intro}

Following Straus \cite{Straus66} and Erd\H{o}s--Nicolas--S\'ark\H{o}zy
\cite{ENS91}, call a set $A$ of positive integers \emph{admissible} if any
two finite subsets of $A$ with the same sum have the same cardinality;
equivalently, the $r$-fold sums of distinct elements of $A$ are pairwise
disjoint for distinct $r$. For finite sets the extremal problem is
essentially closed: intervals $\{N-k+1,\dots,N\}$ with $k\approx2\sqrt N$
are admissible \cite{Straus66}, and $2\sqrt N$ is asymptotically optimal
\cite{DF95,DF99}. For infinite admissible sets, Problem~875 of the
Erd\H{o}s problems collection \cite{EP875} asks how dense such a set can
be, and in particular how small its gaps can be made; Erd\H{o}s remarked
that it is ``not completely trivial'' to find an admissible sequence with
$a_{n+1}/a_n\to1$ (as quoted at \cite{EP875}), a sub-question that remains
open in both directions. Erd\H{o}s--Nicolas--S\'ark\H{o}zy \cite{ENS91}
constructed an infinite admissible set of polynomial counting growth. In
May 2026, a note of GPT-5.5 Pro and Mazur \cite{Mazur26}, with an
accompanying Lean formalization, gave an explicit admissible set with
\[
a_{n+1}-a_n\;\le\; n^{\,3+2\sqrt2}\qquad(3+2\sqrt2=5.828\ldots),
\]
the strongest known gap bound. Neither construction has
$a_{n+1}/a_n\to1$, and \cite{Mazur26} explicitly disclaims the ratio
question.

\medskip
\noindent\textbf{Scope.} \emph{This note resolves none of the questions of
Problem 875.} We prove no new construction, no improvement of the exponent
$3+2\sqrt2$, and nothing about the ratio question. What we do is sharpen
the local machinery underlying \cite{Mazur26} to its exact limits, and
prove two impossibility results that delimit the ``scalar carry''
framework that machinery lives in. Concretely (definitions in
Section~\ref{sec:notation}):

\begin{enumerate}
\item \textbf{Sharp local carry lemma} (Proposition~\ref{prop:sharp}). For
the arithmetic-progression block $C_m=\{1+i(m+1):0\le i<m\}$ and radix
$Q_m=m^2+m+1$, the local carry lemma holds at carry radius exactly
$\lceil m/2\rceil$, with no hypothesis on the quotient, and fails at radius
$\lceil m/2\rceil+1$. The published lemma \cite{Mazur26} has radius $m/4$
under a quotient window $|w|<m/2$. We temper the ``window removed''
reading: at this radius, quotients realized by $C_m$ obey $2|w|<m$
automatically (Remark~\ref{rem:autow}), so the substantive gains are the
doubled, sharp radius and the consequence in item 2.
\item \textbf{Containment-free closure} (Proposition~\ref{prop:closure}).
The closure step of the construction --- adjoining a scaled block $MC_k$ to
an already-good pair $(A,M)$ --- needs only the cardinality bound
$|A|\le\lceil k/2\rceil$. The hypotheses $A\subset[1,M)$ and
$\SSum(A)/M<k/4$ of \cite[Lemma~3]{Mazur26} are dropped entirely, and no
quotient bound is tracked. Block length $\approx2|A|$ suffices in place of
$\approx4|A|$.
\item \textbf{Quadratic radix barrier} (Proposition~\ref{prop:wall}).
Every $k$-element block $C$ satisfying the exact scalar carry property
\emph{at all quotients} has radix $Q\ge\lfloor k^2/4\rfloor+1$ (and
$\ge\lfloor(k+1)^2/4\rfloor+1$ if $Q\notin C$); both bounds are attained,
by flat blocks. So no all-quotient scalar block can beat the quadratic
radix of \cite{Mazur26} in exponent. We stress the qualifier: the barrier
says nothing about invariants imposed only for a \emph{bounded window} of
quotients, which are explicitly not ruled out
(Remark~\ref{rem:window-scope}).
\item \textbf{Rigidity of staggered lifts} (Proposition~\ref{prop:rigid}).
Lifting the Mazur block by multiples of the radix,
$c_i\mapsto c_i+Q_m\lambda_i$, preserves the exact scalar invariant if and
only if $\lambda_i=\lambda(m-i)$ for some $\lambda\in\Z$; carry radius $1$
then forces $\lambda=0$. The rigidity persists for bounded-window
invariants once the window exceeds the lift size
(Remark~\ref{rem:windowed-rigidity}).
\end{enumerate}

Items 1--2 are the positive half: they show the bookkeeping side of the
carry method is not the obstruction to denser or overlapping block
towers, partially meeting the desideratum of \cite[Remark~4]{Mazur26}.
Items 3--4 are the negative half: within the exact, all-quotient, scalar
framework, the radix cost is intrinsically quadratic and the known residue
system cannot be geometrically rearranged. Section~\ref{sec:consequences}
spells out what this leaves open for the two questions of Problem 875.

All proofs are elementary and complete. Propositions \ref{prop:sharp}
and \ref{prop:closure} have additionally been formalized in Lean~4 ---
both directions of Proposition~\ref{prop:sharp}, with the sharpness
proved for every $m$, and zero \texttt{sorry}s; see the formalization
paragraph in Section~\ref{sec:provenance} --- and
Proposition~\ref{prop:sharp}'s frontier has been machine-checked for
$3\le m\le26$ (Remark~\ref{rem:machine}). The mathematics in this note
was produced and cross-verified by AI systems; see the provenance
statement in Section~\ref{sec:provenance}.

\section{Notation and the carry framework of the Mazur
construction}\label{sec:notation}

For a finite set $B\subset\Z$ and $t\in\Z$, the \emph{graded difference
set} is
\[
\Delta_t(B)\;=\;\Bigl\{\sum_{x\in U}x-\sum_{x\in V}x \;:\;
U,V\subseteq B,\ |U|-|V|=t\Bigr\}.
\]
For an infinite set the same definition is used with $U,V$ finite. Thus
$A$ is admissible if and only if $0\notin\Delta_t(A)$ for every $t\ne0$.
Replacing $(U,V)$ by $(U\setminus V,\,V\setminus U)$ changes neither the
difference nor the cardinality defect, so every element of $\Delta_t(B)$,
$B=\{b_1,\dots,b_k\}$, can be written as $\sum_i\varepsilon_ib_i$ with
$\varepsilon\in\{-1,0,1\}^k$ and $\sum_i\varepsilon_i=t$, and conversely.

\begin{definition}\label{def:local}
Let $C\subset\Z$ be finite, $Q\ge1$, $U\ge0$. Say $\loc{U}(C,Q)$ holds if
for all $t,w,u\in\Z$ with $|u|\le U$:
\[
Y\in\Delta_t(C),\quad Y=Qw-u
\qquad\Longrightarrow\qquad t=w-u .
\]
There is no restriction on $w$; we call this the \emph{all-quotient}
property, and $U$ the \emph{carry radius}. Clearly
$\loc{U}\Rightarrow\loc{U'}$ for $U'\le U$.
\end{definition}

\begin{definition}\label{def:wcg}
A pair $(A,M)$, with $A\subset\Z$ finite and $M\ge1$, is \emph{weakly
carry-good} if for all $T,q\in\Z$: $D\in\Delta_T(A)$ and $D=Mq$ imply
$T=q$. Equivalently, $(A,M)$ is weakly carry-good iff $\loc{0}(A,M)$
holds. Taking $D=0=M\cdot0$ shows every weakly carry-good $A$ is
admissible.
\end{definition}

The construction of \cite{Mazur26} uses the blocks and radices
\[
C_m=\{1+i(m+1):0\le i<m\},\qquad Q_m=m^2+m+1 ,
\]
so $|C_m|=m$, $\max C_m=m^2$, $\SSum(C_m)=m(m^2+1)/2$, and
$Q_m-\max C_m=m+1$ equals the internal spacing. Every element of $C_m$ is
$\equiv1\pmod{m+1}$, so a sum of $r$ distinct elements is $\equiv
r\pmod{m+1}$: the block is admissible with the residue reading off the
cardinality. Globally, blocks $B_j=M_jC_{k_j}$ are stacked with moduli
$M_{j+1}=M_jQ_{k_j}$ and lengths $k_j\approx2(1+N_{j-1})^{1+\sqrt2}$
($N_{j-1}$ = number of prior elements); the blocks occupy consecutive
intervals $[M_j,\,M_jk_j^2]$ and all gaps at stage $j$ equal
$M_j(k_j+1)$, which yields the exponent $3+2\sqrt2$. Two lemmas drive the
induction \cite{Mazur26}:

\begin{itemize}
\item \emph{Local carry lemma} (their Lemma 2): if $Y\in\Delta_t(C_m)$ and
$Y=Q_mw-u$ with $4|u|<m$ and $2|w|<m$, then $t=w-u$. (These are the exact
hypotheses of the Lean artifact \texttt{LocalCarry.lean}.)
\item \emph{Carry closure} (their Lemma 3): call $(A,M)$ \emph{carry-good}
if $A\subset[1,M)\cap\N$ and the implication of Definition~\ref{def:wcg}
holds. If $(A,M)$ is carry-good and $\SSum(A)/M<k/4$, then
$(A\cup MC_k,\,MQ_k)$ is carry-good, with $\SSum(A')/M'<k/2$ (which is
what bounds the quotient $|w|$ fed to Lemma 2 at the next stage).
\end{itemize}

Sections \ref{sec:sharp}--\ref{sec:closure} replace both lemmas by
sharper, hypothesis-lighter statements; Sections
\ref{sec:wall}--\ref{sec:rigidity} bound what any such improvement can
achieve.

\section{The sharp local carry lemma}\label{sec:sharp}

\begin{proposition}[sharp local carry lemma]\label{prop:sharp}
For every $m\ge1$, $\loc{\lceil m/2\rceil}(C_m,Q_m)$ holds. This is sharp:
$\loc{\lceil m/2\rceil+1}(C_m,Q_m)$ fails.
\end{proposition}

\begin{proof}
Write $d=m+1$, $Q=Q_m=md+1$, $L=m(m+1)/2$, and $r=\lceil m/2\rceil$. The
$i$-th element of $C_m$ is
\[
c_i \;=\; 1+di \;=\; Q-d(m-i)\qquad(0\le i<m).
\]

\emph{The lemma at radius $r$.} Let $Y\in\Delta_t(C_m)$; as in
Section~\ref{sec:notation} write $Y=\sum_{i=0}^{m-1}\varepsilon_ic_i$ with
$\varepsilon_i\in\{-1,0,1\}$ and $t=\sum_i\varepsilon_i$. Put
$z=\sum_i(m-i)\varepsilon_i$; then
\[
Y \;=\; Qt-dz,
\qquad
|z|\;\le\;\sum_{i=0}^{m-1}(m-i)\;=\;L .
\]
Suppose $Y=Qw-u$ with $|u|\le r$, and set $v=w-t$, so that $Qv=u-dz$.
Since $Q\equiv1\pmod d$, reducing modulo $d$ gives $v\equiv u\pmod d$,
say $v=u+\ell d$ with $\ell\in\Z$. Substituting back and using $Q-1=md$,
\begin{equation}\label{eq:key}
z \;=\; -mu-\ell Q .
\end{equation}
If $\ell\ne0$ then
\[
|z| \;=\; |mu+\ell Q| \;\ge\; |\ell|\,Q-m|u| \;\ge\; Q-mr \;>\; L,
\]
where the last inequality is a parity check:
\[
Q-mr-L=
\begin{cases}
h+1, & m=2h,\\
1, & m=2h+1 .
\end{cases}
\]
This contradicts $|z|\le L$. Hence $\ell=0$, so $v=u$, i.e.\
$t=w-u$.\footnote{The Lean formalization (Section~\ref{sec:provenance})
surfaced a simplification: no parity split is actually needed here, since
$2mr\le m(m+1)=2L$ and $Q=2L+1$ give $Q-mr-L\ge1$ in both parities at
once. We keep the explicit split, which the sharpness direction uses
anyway.}

\emph{Sharpness.} Let $u_0=r+1$ and $z_0=Q-mu_0$. Explicitly,
\begin{align*}
m=2h\ (h\ge1):&\qquad z_0=2h^2+1\;\le\;L=2h^2+h,\\
m=2h+1\ (h\ge0):&\qquad z_0=2h^2+h+1\;\le\;L=2h^2+3h+1,
\end{align*}
so $0\le z_0\le L$ in both parities. The coefficients $m-i$
($0\le i<m$) are exactly $1,2,\dots,m$, and every integer in $[0,L]$ is a
sum of a subset of $\{1,\dots,m\}$ (induction on $m$: if $s\ge m$ use $m$
and apply the inductive step to $s-m\le(m-1)m/2$; if $s<m$ then $s$ is
itself an element or $0$). Choose $\varepsilon\in\{0,1\}^m$ with
$\sum_i(m-i)\varepsilon_i=z_0$, and put $t=\sum_i\varepsilon_i$, so
$Y=Qt-dz_0\in\Delta_t(C_m)$. With $w=t+u_0-d$ we get, using $Q-1=md$ and
$z_0=Q-mu_0$,
\[
Qw-u_0 \;=\; Qt+(Q-1)u_0-Qd \;=\; Qt+d(mu_0-Q) \;=\; Qt-dz_0 \;=\; Y,
\]
while $w-u_0=t-d\ne t$. Thus $\loc{r+1}(C_m,Q_m)$ fails.
\end{proof}

\begin{remark}[machine check]\label{rem:machine}
Two independent implementations corroborate the frontier: exhaustive
enumeration over all $\varepsilon\in\{-1,0,1\}^m$ for $3\le m\le24$ (log
in the companion repository:\footnote{File paths in this note refer to
the companion repository,
\url{https://github.com/Mister-Meeseeks/erdos-875-carry-lemmas},
which also contains the scripts reproducing both checks.}
\path{verification/machine-check-prop1-frontier-m24.txt}) and
a dynamic program over all graded differences for $3\le m\le26$
(log: \path{verification/machine-check-frontier-m26.txt}) both find the
minimal $|u|$ admitting a violation to be exactly $\lceil m/2\rceil+1$
in every case; $m\le2$ was checked by hand. For
example, at $m=3$ the first violation is $Y=10=c_0+c_2=13\cdot1-3$, with
$t=2\ne w-u=-2$ --- this is the witness produced by the proof above
($u_0=3$, $z_0=4=3+1$, $\varepsilon=(1,0,1)$, $w=1$). These checks are
not load-bearing --- the proof above covers every $m$, and both
directions are also machine-verified in Lean for all $m$
(Section~\ref{sec:provenance}) --- they corroborate the frontier by
independent means.
\end{remark}

\begin{remark}[how much of the quotient window was real]\label{rem:autow}
Proposition~\ref{prop:sharp} carries no hypothesis on $w$, but for the
block $C_m$ this is largely automatic rather than extra strength: if
$Y\in\Delta_t(C_m)$ and $Y=Q_mw-u$ with $|u|\le m$, then
$Q_m|w|\le|Y|+|u|\le\SSum(C_m)+m$ with $\SSum(C_m)=m(m^2+1)/2$, so
\[
|w|\;\le\;\Bigl\lfloor\frac{\SSum(C_m)+m}{Q_m}\Bigr\rfloor
\;\le\;\Bigl\lfloor\frac{m}{2}\Bigr\rfloor
\]
for every \emph{realizable} quotient (the first floor can be strictly
smaller: at $m=4$ it is $1$, not $\lfloor m/2\rfloor=2$). Sharper: at the
radius of Proposition~\ref{prop:sharp}, $|u|\le\lceil m/2\rceil$, the
same computation gives $|w|\le\lfloor(m-1)/2\rfloor$, i.e.\ $2|w|<m$
automatically --- every realizable quotient lies strictly \emph{inside}
the strict window $2|w|<m$ of \cite[Lemma~2]{Mazur26}, so that window was
never a live constraint on realizable representations. The substantive
content of Proposition~\ref{prop:sharp} is the sharp radius
$\lceil m/2\rceil$ (versus $m/4$), together with the fact that the
all-quotient formulation lets the closure step below dispense with any
quotient accounting: there the quotient is an a priori unbounded divisor
quotient, not yet known to be realizable.
\end{remark}

\section{Containment-free closure}\label{sec:closure}

\begin{proposition}[containment-free closure]\label{prop:closure}
Let $(A,M)$ be weakly carry-good (Definition~\ref{def:wcg}), let $N=|A|$,
let $k\ge1$, and suppose
\[
N\;\le\;\Bigl\lceil\frac{k}{2}\Bigr\rceil,
\qquad
A\cap MC_k=\varnothing .
\]
Then $\bigl(A\cup MC_k,\;MQ_k\bigr)$ is weakly carry-good.
\end{proposition}

\begin{proof}
Let $D\in\Delta_T(A\cup MC_k)$ with $D=MQ_kq$; we must show $T=q$. Since
$A$ and $MC_k$ are disjoint, any pair of subsets of the union splits, so
\[
D=D_A+MY,\qquad D_A\in\Delta_{T_A}(A),\quad Y\in\Delta_t(C_k),\quad
T=T_A+t .
\]
From $D=MQ_kq$ we get $D_A=M(Q_kq-Y)=Me$ with $e:=Q_kq-Y\in\Z$. Weak
carry-goodness of $(A,M)$ gives $e=T_A$, whence
\[
|e|=|T_A|\le|A|=N\le\Bigl\lceil\frac{k}{2}\Bigr\rceil .
\]
Now $Y=Q_kq-e$ is a representation with carry $|e|\le\lceil k/2\rceil$
and quotient $q$; Proposition~\ref{prop:sharp} (all-quotient, radius
$\lceil k/2\rceil$) yields $t=q-e$. Therefore
$T=T_A+t=e+(q-e)=q$.
\end{proof}

\begin{remark}[exactly which hypotheses are dropped]\label{rem:dropped}
Proposition~\ref{prop:closure} is the direct upgrade of
\cite[Lemma~3]{Mazur26}. Compared with that statement it drops:
\begin{enumerate}
\item the containment $A\subset[1,M)\cap\N$ --- here $A$ is an arbitrary
finite set of integers, which is the shape of hypothesis an
\emph{overlapping}-scales construction would need;
\item the mass estimate $\SSum(A)/M<k/4$ --- replaced by the cardinality
bound $N\le\lceil k/2\rceil$, i.e.\ $k\ge2N-1$, so blocks of length
$\approx2|A|$ suffice where $\approx4|A|$ was needed;
\item all quotient bookkeeping --- no analogue of the output estimate
$\SSum(A')/M'<k/2$ is tracked, because Proposition~\ref{prop:sharp} needs
no quotient window. This is where the all-quotient form earns its keep:
in the proof, $q$ is a divisor quotient not a priori bounded by anything.
\end{enumerate}
The conclusion is again weak carry-goodness, so the statement iterates
and recovers the original induction (carry-good $\Rightarrow$ weakly
carry-good).
\end{remark}

\begin{remark}[what is \emph{not} improved]\label{rem:geometry}
Proposition~\ref{prop:closure} does not change the geometry of the
construction: for the standard block, $\max C_k=k^2<Q_k$, so the tower
$M_{j+1}=M_jQ_{k_j}$ still stacks blocks in disjoint consecutive
intervals, with the same gaps and the same exponent $3+2\sqrt2$. What the
proposition shows is that the \emph{closure mechanism} imposes no
containment or mass obstruction on more ambitious (e.g.\ overlapping)
geometries --- only the local lemma and the block geometry do. This
partially meets the desideratum of \cite[Remark~4]{Mazur26}, which asks
what a replacement local block must supply: by
Proposition~\ref{prop:closure}, only a carry lemma of radius
$\ge|A|$ compatible with its own radix, plus the geometric wish-list.
\end{remark}

\section{A quadratic radix barrier for all-quotient scalar
carry}\label{sec:wall}

The results above sharpen the known block; the next one bounds every
possible block within the exact scalar framework. Throughout this
section the \emph{all-quotient} qualifier is essential; see
Remark~\ref{rem:window-scope}.

For a finite $D\subset\Z$ let $\FS(D)=\{\sum_{x\in S}x:S\subseteq D\}$
denote the set of distinct subset sums, including the empty sum $0$.

\begin{lemma}[carry $\Rightarrow$ modular injectivity]\label{lem:inj}
Let $C\subset\Z$ with $|C|=k$ and let $Q\ge1$ satisfy $\loc0(C,Q)$, i.e.
\begin{equation}\label{eq:exact}
Y\in\Delta_t(C),\quad Y=Qw
\qquad\Longrightarrow\qquad t=w .
\end{equation}
Then reduction modulo $Q$ is injective on $\FS(C-Q)$; in particular
$|\FS(C-Q)|\le Q$.
\end{lemma}

\begin{proof}
Put $D=C-Q=\{c-Q:c\in C\}$; subsets of $D$ correspond bijectively to
subsets $X\subseteq C$, with $\sigma_D(X)=\sigma_C(X)-Q|X|$, where
$\sigma$ denotes the subset sum. Suppose
$\sigma_D(X)\equiv\sigma_D(X')\pmod Q$ for $X,X'\subseteq C$. Then
$\sigma_C(X)-\sigma_C(X')\equiv0\pmod Q$, say
$\sigma_C(X)-\sigma_C(X')=Qw$; this difference lies in $\Delta_t(C)$ with
$t=|X|-|X'|$, so \eqref{eq:exact} forces $w=t$. Hence
\[
\sigma_D(X)-\sigma_D(X')
=\bigl(\sigma_C(X)-\sigma_C(X')\bigr)-Q\bigl(|X|-|X'|\bigr)
=Qw-Qt=0 .
\]
So congruent values of $\FS(D)$ are equal, and $\FS(D)$ injects into
$\Z/Q\Z$.
\end{proof}

The counting input we need is known: the following is the $\alpha=0$ case
of \cite[Corollary~9]{BP22} (proved there via their Theorems 4 and 7),
and it is optimal in that case by \cite[Remark~10]{BP22}; the $\alpha=1$
analogue goes back to Nathanson \cite[Theorem~3]{Nathanson95}. We include
an elementary proof for completeness, so that this note is
self-contained.

\begin{lemma}[subset-sum counting; known]\label{lem:count}
Let $D$ be a set of $k$ distinct integers. Then
\[
|\FS(D)|\;\ge\;\Bigl\lfloor\frac{k^2}{4}\Bigr\rfloor+1,
\]
and if $0\notin D$ then
$|\FS(D)|\ge\bigl\lfloor(k+1)^2/4\bigr\rfloor+1$.
\end{lemma}

\begin{proof}
Let $p$ and $n$ be the numbers of positive and negative elements of $D$,
so $p+n=s:=k-[0\in D]$.

If $b_1<\dots<b_p$ are the positive elements, then the sums
\[
T_j+b_i,\qquad 1\le j\le p,\quad 1\le i\le p-j+1,
\qquad T_j:=b_{p-j+2}+\dots+b_p\ \ (T_1=0),
\]
form a strictly increasing chain of $p+(p-1)+\dots+1=p(p+1)/2$ subset
sums: each row is increasing in $i$, the last entry of row $j$ is
$T_{j+1}$, and the first entry of row $j+1$ exceeds $T_{j+1}$ because
$b_1>0$. So $D$ has at least $p(p+1)/2$ distinct \emph{positive} subset
sums, and symmetrically at least $n(n+1)/2$ distinct \emph{negative}
ones; together with the empty sum $0$ these are pairwise distinct, giving
\[
|\FS(D)|\;\ge\;\frac{p(p+1)}{2}+\frac{n(n+1)}{2}+1 .
\]
For fixed $s=p+n$ the right-hand side, as a function of $p$, changes by
$2p+1-s$ when $p\mapsto p+1$, so it is minimized at the balanced split,
where its value is $\lfloor(s+1)^2/4\rfloor+1$ (both parities checked
directly). If $0\notin D$ then $s=k$ and we get
$\lfloor(k+1)^2/4\rfloor+1$; if $0\in D$ then $s=k-1$ and we get
$\lfloor k^2/4\rfloor+1$. Since
$\lfloor(k+1)^2/4\rfloor\ge\lfloor k^2/4\rfloor$, the first display holds
in both cases.
\end{proof}

\begin{proposition}[quadratic radix barrier]\label{prop:wall}
Let $C$ be a set of $k$ distinct integers and let $Q\ge1$ satisfy the
exact \emph{all-quotient} scalar carry property \eqref{eq:exact} ---
equivalently, let $(C,Q)$ be weakly carry-good. Then
\[
Q\;\ge\;\Bigl\lfloor\frac{k^2}{4}\Bigr\rfloor+1 ,
\]
and if moreover $Q\notin C$ then
$Q\ge\bigl\lfloor(k+1)^2/4\bigr\rfloor+1$. Both bounds are attained for
every $k$ (Remark~\ref{rem:wall-sharp}).
\end{proposition}

\begin{proof}
The set $D=C-Q$ has $k$ distinct elements, and $0\in D$ iff $Q\in C$.
Combine Lemma~\ref{lem:inj} with Lemma~\ref{lem:count}.
\end{proof}

Consequently any block family retaining the exact all-quotient scalar
invariant has $Q=\Omega(k^2)$: the quadratic radix of the Mazur block
($Q_k=k^2+k+1$) is optimal \emph{in exponent} for the entire all-quotient
scalar-carry framework, not merely for arithmetic progressions. The
hoped-for all-quotient scalar block with $Q\asymp k^{2-\delta}$ does not
exist.
This does not prove the gap exponent $3+2\sqrt2$ optimal among all
constructions; see Section~\ref{sec:consequences}.

\begin{remark}[scope: the barrier does \emph{not} touch bounded-window
invariants]\label{rem:window-scope}
Every statement of the barrier above carries the all-quotient qualifier,
and it cannot be removed: the proof of Lemma~\ref{lem:inj} applies
\eqref{eq:exact} to differences $\sigma_C(X)-\sigma_C(X')$ whose
quotients $w$ can be as large as $\bigl(\sum_{c\in C}|c|\bigr)/Q$ --- and
$\sum_{c\in C}|c|\ge\lfloor k^2/4\rfloor$ for any $k$-element set, so in
the regime being excluded ($Q$ subquadratic in $k$) these quotients are
\emph{unbounded} in $k$. A local property required only for
$|w|\le W$ with, say, $W\approx\varepsilon m^2$ --- the ``wide-window''
regime that a multiplicatively overlapping tower would actually need ---
is untouched by Proposition~\ref{prop:wall}. Bounded-window invariants
with subquadratic radix are \emph{not} ruled out by anything in this
note, and finding them (over new residue systems, cf.\
Proposition~\ref{prop:rigid}) is in our view the live attack direction on
both questions of Problem 875.
\end{remark}

\begin{remark}[the counting step is tight, and the barrier is
sharp]\label{rem:wall-sharp}
Optimality at $\alpha=0$ in \cite[Remark~10]{BP22} means the counting
step (Lemma~\ref{lem:count}) cannot be improved, so any improvement of
the barrier would have to strengthen the injectivity step. In fact none
is possible: the barrier itself is sharp. If
$\sum_{c\in C}|c-Q|<Q$, then $(C,Q)$ satisfies \eqref{eq:exact}: writing
$Y=\sum_i\varepsilon_ic_i=Qt+\sum_i\varepsilon_i(c_i-Q)$, the second
term has absolute value $<Q$, so $Y=Qw$ forces $w=t$. Such ``flat''
blocks are exactly the interval-separation regime of Straus
\cite{Straus66} and \cite{ENS91}. Now the minimum of $\sum_{d\in D}|d|$ over $k$-element
integer sets $D$ is $\lfloor k^2/4\rfloor$ (balanced interval around $0$,
which contains $0$), and $\lfloor(k+1)^2/4\rfloor$ if $0\notin D$ is
required. Taking $C=Q+D$ with $D$ extremal and
$Q=\lfloor k^2/4\rfloor+1$ (resp.\ $\lfloor(k+1)^2/4\rfloor+1$) gives
equality in Proposition~\ref{prop:wall} for every $k$, in both the
$Q\in C$ and $Q\notin C$ cases. Explicit extremizers in both parities
(verified in the targeted critique run cited in
Section~\ref{sec:provenance}): for $k=2h$ take
$D=\{-(h-1),\dots,h\}$ ($0\in D$) or $D=\{-h,\dots,-1,1,\dots,h\}$
($0\notin D$); for $k=2h+1$ take $D=\{-h,\dots,h\}$ or
$D=\{-h,\dots,-1,1,\dots,h+1\}$; in every case $C=Q+D$ consists of
distinct positive integers. At $k=2$ the sharpened $Q\notin C$ bound is
attained by $C=\{2,4\}$, $Q=3$, and the plain bound by $C=\{2,3\}$,
$Q=2$; the degenerate case $k=1$ is attained by $C=\{1\}$, $Q=2$.
Thus the barrier identifies the flat block as the radix-cheapest scalar
block; the \emph{spanning} block of \cite{Mazur26}, with
$\min C_k=1$ and $\max C_k=k^2$, pays only a factor $4+o(1)$ more.
Whether a spanning all-quotient scalar block can have
$Q\le(1-\delta)k^2$ we leave open; Proposition~\ref{prop:wall} caps what
is achievable at $k^2/4$ regardless.
\end{remark}

\begin{remark}[attribution of the technique]\label{rem:technique}
The proof scheme of Proposition~\ref{prop:wall} --- an injectivity
statement modulo $Q$ combined with a quadratic subset-sum count --- is
standard in neighbouring problems, and we make no claim of methodological
novelty: quadratic lower bounds for subset sums in $\Z_p$ go back to
Olson \cite{Olson68} (see also Erd\H{o}s--Heilbronn), the general
abelian-group version is DeVos--Goddyn--Mohar--\v{S}\'amal \cite{DGMS07}
(whose theorem, applied inside $\Z/Q\Z$, would already give
$Q=\Omega(k^2)$ here with a weaker constant), and the
pigeonhole ``distinct subset sums mod $Q$ forces $Q\ge|\FS|$'' is the
engine of the modular Erd\H{o}s distinct-subset-sums problem
\cite{CGKL23}. What is new here is only the reduction
(Lemma~\ref{lem:inj}) from the carry property to modular injectivity,
and the application: a lower bound on the radix of a carry-compatible
block, a quantity \cite[Remark~4]{Mazur26} leaves unconstrained.
\end{remark}

\section{Rigidity of staggered lifts of the Mazur block}\label{sec:rigidity}

A natural attempt to make block towers overlap is to keep the Mazur
residues modulo $Q=Q_m$ but push different elements up by multiples of
$Q$: replace $c_i$ by $c_i+Q\lambda_i$. First, an immediate obstruction
to cruder variants:

\begin{observation}[no repeated residues]\label{obs:dup}
If $(C,Q)$ satisfies \eqref{eq:exact}, then no two distinct elements of
$C$ are congruent modulo $Q$: their difference would be $Qw$ with
$w\ne0$ and cardinality defect $t=0$, forcing $t=w$, a contradiction. In
particular a union of two or more $Q$-translated copies of a block never
satisfies the exact all-quotient scalar invariant.
\end{observation}

Staggered lifts avoid this by using each residue once. They fail anyway:

\begin{proposition}[rigidity of staggered lifts]\label{prop:rigid}
Let $m\ge4$, $Q=Q_m$, $d=m+1$, and for
$\lambda=(\lambda_0,\dots,\lambda_{m-1})\in\Z^m$ set
\[
\widetilde C_\lambda=\{\,c_i+Q\lambda_i \;:\; 0\le i<m\,\},
\qquad c_i=1+di .
\]
\begin{enumerate}
\item[(i)] The $c_i$ are pairwise incongruent modulo $Q$; hence
$|\widetilde C_\lambda|=m$ for every $\lambda$, and distinct lifted
elements never collide.
\item[(ii)] $\widetilde C_\lambda$ satisfies
$\loc0(\widetilde C_\lambda,Q)$ if and only if there is $\lambda\in\Z$
with
\[
\lambda_i=\lambda\,(m-i)\qquad(0\le i<m).
\]
\item[(iii)] If $\widetilde C_\lambda$ satisfies
$\loc1(\widetilde C_\lambda,Q)$, then $\lambda_i=0$ for all $i$. (The
trivial lift satisfies $\loc{\lceil m/2\rceil}$ by
Proposition~\ref{prop:sharp}, so the survivor of (ii) at radius $1$ is
exactly the original block.)
\end{enumerate}
\end{proposition}

\begin{proof}
(i) The $c_i$ are distinct and lie in $[1,m^2]\subset[0,Q)$.

(ii), necessity. For $1\le j\le m-2$ the three indices $j-1,j,m-1$ are
distinct and
\[
c_j+c_{m-1}-c_{j-1} \;=\; 1+dm \;=\; Q ,
\]
an element of $\Delta_1$ with quotient $1$. After lifting,
\[
\widetilde c_j+\widetilde c_{m-1}-\widetilde c_{j-1}
= Q\bigl(1+\lambda_j+\lambda_{m-1}-\lambda_{j-1}\bigr)\in
\Delta_1(\widetilde C_\lambda),
\]
so $\loc0$ forces $1=1+\lambda_j+\lambda_{m-1}-\lambda_{j-1}$, i.e., with
$\lambda:=\lambda_{m-1}$,
\begin{equation}\label{eq:affine}
\lambda_{j-1}-\lambda_j=\lambda\qquad(1\le j\le m-2),
\end{equation}
whence $\lambda_j=\lambda_0-j\lambda$ for $0\le j\le m-2$. It remains to
pin the intercept $\lambda_0=m\lambda$.

If $m\ge5$, the indices $0,2,m-2$ are distinct and
$c_2+c_{m-2}-c_0=1+dm=Q$, so $\loc0$ forces
$\lambda_2+\lambda_{m-2}-\lambda_0=0$; substituting
$\lambda_2=\lambda_0-2\lambda$ and $\lambda_{m-2}=\lambda_0-(m-2)\lambda$
gives $\lambda_0=m\lambda$.

If $m=4$ (where $2=m-2$ makes the previous relation degenerate), use
instead $c_1+c_2-c_0-c_3=0\in\Delta_0$ with quotient $0$: $\loc0$ forces
$\lambda_1+\lambda_2-\lambda_0-\lambda_3=0$, and with
$\lambda_1=\lambda_0-\lambda$, $\lambda_2=\lambda_0-2\lambda$,
$\lambda_3=\lambda$ this gives $\lambda_0=4\lambda$.

In both cases $\lambda_j=(m-j)\lambda$ for $0\le j\le m-2$, and
$\lambda_{m-1}=\lambda=(m-(m-1))\lambda$.

(ii), sufficiency. Suppose $\lambda_i=\lambda(m-i)$. Any
$\widetilde Y\in\Delta_t(\widetilde C_\lambda)$ has the form
$\widetilde Y=\sum_i\varepsilon_i\widetilde c_i$ with
$\varepsilon\in\{-1,0,1\}^m$, $t=\sum_i\varepsilon_i$; putting
$z=\sum_i(m-i)\varepsilon_i$ as before,
\[
\widetilde Y \;=\; (Qt-dz)+Q\lambda z \;=\; Qt+(Q\lambda-d)z,
\qquad |z|\le L=\tfrac{m(m+1)}2 .
\]
If $\widetilde Y=Qw$, then $Q\mid dz$; since $Q=md+1$ is coprime to $d$,
$Q\mid z$, and $|z|\le L<Q$ forces $z=0$. Hence $\widetilde Y=Qt$ and
$w=t$, which is $\loc0$.

(iii) $\loc1$ implies $\loc0$, so by (ii) $\lambda_i=\lambda(m-i)$. The
negative singleton gives
\[
-\widetilde c_0=-1-Q\lambda_0=Q(-\lambda_0)-1\in
\Delta_{-1}(\widetilde C_\lambda),
\]
i.e.\ $t=-1$, $w=-\lambda_0$, $u=1$. $\loc1$ forces $-1=-\lambda_0-1$,
so $\lambda_0=0$; by (ii) $\lambda_0=m\lambda$, hence $\lambda=0$ and all
$\lambda_i$ vanish.
\end{proof}

\begin{remark}[the rigidity persists for bounded
windows]\label{rem:windowed-rigidity}
The proof of necessity in (ii) and of (iii) evaluates the invariant only
at the quotients
$w=1+\lambda_j+\lambda_{m-1}-\lambda_{j-1}$,
$w=\lambda_1+\lambda_2-\lambda_0-\lambda_3$ (for $m=4$),
$w=1+\lambda_2+\lambda_{m-2}-\lambda_0$, and $w=-\lambda_0$, all of
absolute value at most $4\|\lambda\|_\infty+1$, and only at carry radii
$u\in\{0,1\}$. Hence the conclusions of
Proposition~\ref{prop:rigid}(ii)--(iii) already follow from the
\emph{windowed} invariant
\[
Y\in\Delta_t(\widetilde C_\lambda),\ \ Y=Qw-u,\ \ |u|\le1,\ \
|w|\le 4\|\lambda\|_\infty+1
\ \Longrightarrow\ t=w-u .
\]
So staggered lifts of the Mazur residue system are dead \emph{even in the
bounded-window regime} of Remark~\ref{rem:window-scope}, as soon as the
window is comparable to the lift size --- window growth cannot be bought
by re-using these residues. A wide-window construction must use a
genuinely new residue system.
\end{remark}

\begin{remark}[the threshold $m\ge4$ is sharp]\label{rem:m3}
For $m=3$ ($C_3=\{1,5,9\}$, $Q=13$) the only relations available to
$\loc0$ are $\pm(c_1+c_2-c_0)=\pm Q$, and one checks that
$\loc0(\widetilde C_\lambda,13)$ holds precisely when
$\lambda_0=\lambda_1+\lambda_2$ --- a two-parameter family strictly
containing $\{\lambda(3,2,1):\lambda\in\Z\}$. So $\lambda=(1,1,0)$
survives $\loc0$: this establishes sharpness of the threshold for the
exact classification of Proposition~\ref{prop:rigid}(ii) specifically.
The threshold is also necessary for part (iii): $\lambda=(0,s,-s)$
satisfies $\loc1(\widetilde C_\lambda,13)$ for every $s\in\Z$, because
every signed relation at $m=3$ with $|u|\le1$ has
$\varepsilon_1=\varepsilon_2$, so the lift contribution
$13s(\varepsilon_1-\varepsilon_2)$ vanishes (witness from the targeted
critique run cited in Section~\ref{sec:provenance}).
\end{remark}

\begin{remark}[context]\label{rem:mr}
The closest prior result in spirit is Monopoli--Ruzsa \cite{MR15}
(following work of Diaconis, Shao and Soundararajan on carries): digit
sets modulo $m^2$ inducing minimal carry behaviour under two-summand
addition are forced to be arithmetic progressions up to linear maps ---
an ``AP structure is rigid under carry constraints'' theorem. The setting
there (two-summand addition carries, no cardinality grading, no lifts of
a fixed residue system, no admissibility) is disjoint from ours, and no
statement transfers in either direction; we cite it as the moral
precedent for Proposition~\ref{prop:rigid}.
\end{remark}

\section{Consequences for Problem 875}\label{sec:consequences}

Nothing in this note answers either question of \cite{EP875}. What it
does is redistribute where the difficulty provably lives.

\paragraph{The gap exponent.} Redoing the block-schedule optimization of
\cite[Remark~2]{Mazur26} with a hypothetical local block of $k$ elements
and radix $Q\asymp k^\gamma$ (heuristic scaling, assuming the same
closure mechanics: the exponent condition becomes
$\tau+\theta(\gamma-\tau)\le0$) gives optimized gap exponent
$(1+\sqrt\gamma\,)^2$; the value $\gamma=2$ of the standard block gives
$3+2\sqrt2$, and any all-quotient scalar block with $\gamma<2$ would have
beaten it. Proposition~\ref{prop:wall} closes exactly that route:
$\gamma\ge2$ is forced for every all-quotient scalar block, so within
this framework the exponent $3+2\sqrt2$ is optimal, and by
Remark~\ref{rem:wall-sharp} even the constant has at most a factor
$4+o(1)$ of slack. Improving the exponent therefore requires leaving the
framework, i.e.\ relaxing at least one of:
\begin{enumerate}
\item the \emph{all-quotient} requirement --- a local lemma valid only
for $|w|\le W$ (Remark~\ref{rem:window-scope}); this is untouched by our
barrier and is the direction we consider live;
\item the \emph{scalar} invariant $q=t$ --- vector-valued or multi-digit
invariants;
\item the single-stage closure by one modulus --- genuinely multi-level
local structures;
\item the Mazur residue system itself --- and Proposition~\ref{prop:rigid}
with Remark~\ref{rem:windowed-rigidity} shows any lift-based overlap
attempt is forced to abandon it anyway.
\end{enumerate}

\paragraph{The ratio question.} For $a_{n+1}/a_n\to1$ one needs the
enumeration to meet every multiplicative window, which the
nested-interval tower cannot do at block entry; scales must overlap.
Propositions \ref{prop:sharp} and \ref{prop:closure} clear two suspects:
neither the quotient window nor the containment/mass hypotheses of the
closure were the obstruction --- the closure step now needs only a
cardinality bound, and accepts an arbitrary finite $A\subset\Z$. What
remains is geometry: $\max C_k<Q_k$ keeps the tower non-overlapping
(Remark~\ref{rem:geometry}), Proposition~\ref{prop:rigid} says the known
residues cannot be lifted past the next modulus, and
Proposition~\ref{prop:wall} says the radix cannot be shrunk to close the
entry gap within the all-quotient scalar framework. The surviving routes
are the same four as above, with (1)+(4) --- bounded-window local lemmas
over new residue systems --- the concrete target: one wants residues
$r_1,\dots,r_k\pmod Q$ whose protected relation lattice at radius $U$ and
window $W$ leaves enough lift freedom to cover a long multiplicative
interval, together with a global argument confining all realized
quotients to the window. Whether such finite-window freedom survives the
global induction is open; Proposition~\ref{prop:wall} only says its
all-quotient idealization does not.

\begin{question}
Is there a block family $(C^{(m)},Q^{(m)})$, $|C^{(m)}|=m$, satisfying
the scalar carry property at radius $U(m)\gtrsim m$ for all quotients
$|w|\le W(m)$ with $W(m)/m\to\infty$ (ideally $W\approx\varepsilon m^2$),
such that either $Q^{(m)}=o(m^2)$ or $\max C^{(m)}\gg Q^{(m)}$ (so that
towers can overlap)? By Remark~\ref{rem:windowed-rigidity} its residue
system cannot be a staggered lift of $C_m$; flat blocks
(Remark~\ref{rem:wall-sharp}) have spread below $Q$ and cannot overlap,
so it must be genuinely spanning.
\end{question}

\section{Provenance}\label{sec:provenance}

This note is a machine-generated artifact and we state its provenance
plainly; nothing here should be read as a claim of autonomous human-free
discovery, nor as more than what is described.

The four propositions and their proofs were produced by an OpenAI model
(\texttt{gpt-5.6-sol}) during a single attack run in the
\textsc{erdosSweep} harness on 2026-07-09 (transcript in the companion
repository:
\path{provenance/attack-transcript-codex.md}).
They were then adversarially verified, line by line, by an Anthropic
model (Claude, Fable~5) in a separate read-only session (transcript:
\path{provenance/critique-transcript-claude.md}),
which re-derived every algebraic step, refuted nothing, flagged two
presentation-level overstatements (the missing all-quotient qualifier
and the ``window removed'' framing, both corrected here), and
contributed several strengthenings incorporated here: the self-contained proof of
Lemma~\ref{lem:count}, the no-collision statement and the converse
(sufficiency) direction in Proposition~\ref{prop:rigid}, the windowed
rigidity of Remark~\ref{rem:windowed-rigidity}, and the scope
corrections now embedded in Remarks \ref{rem:autow} and
\ref{rem:window-scope}. Proposition~\ref{prop:sharp}'s frontier was
machine-checked by two independent implementations for $3\le m\le26$
(Remark~\ref{rem:machine}); $m\le2$ by hand. A literature check
(2026-07-10; report
\path{provenance/literature-check.md})
read the primary sources \cite{Mazur26} (note and Lean repository) and
\cite{BP22} in full and swept the adjacent subset-sum, carries and
admissible-set literature: no prior statement of Propositions
\ref{prop:sharp}--\ref{prop:rigid} was found, while the counting
ingredient of Proposition~\ref{prop:wall} is known and attributed
(Lemma~\ref{lem:count}, Remark~\ref{rem:technique}); MathSciNet and
zbMATH were not accessible, which is the residual risk for
technique-level precedents of Proposition~\ref{prop:wall}.

\emph{Formalization.} Propositions \ref{prop:sharp} and
\ref{prop:closure} are formalized in Lean~4 with mathlib (directory
\path{lean/} in the companion repository; toolchain
\path{leanprover/lean4:v4.30.0-rc2} with the matching mathlib tag, the
same pins as the Lean repository of \cite{Mazur26}; no code was reused
from that repository --- the definitions are written fresh against the
statements here, with $\Delta_t$ encoded via \texttt{Finset} pairs). The
build is green with zero \texttt{sorry}s, and \texttt{\#print axioms} on
the main theorems reports only the standard \texttt{propext},
\texttt{Classical.choice} and \texttt{Quot.sound}. The main declarations
are: \path{Erdos875.local_carry} --- Proposition~\ref{prop:sharp},
main direction, with the radius stated division-free as $2|u|\le m+1$
and \emph{no} hypothesis on the quotient $w$;
\path{Erdos875.local_carry_sharp} --- the sharpness direction, an
explicit violation at $|u|=\lceil m/2\rceil+1$ for \emph{every} $m\ge1$
(a theorem for all $m$, not merely the machine-checked instances); and
\path{Erdos875.weaklyCarryGood_union} --- Proposition~\ref{prop:closure},
with only the cardinality hypothesis $2|A|\le k+1$, disjointness,
$M\ne0$ and $k\ge1$. Propositions \ref{prop:wall} and \ref{prop:rigid}
are \emph{not} formalized. The formalization also surfaced the
simplification recorded in the footnote to
Proposition~\ref{prop:sharp}'s proof (the parity split in the main
direction is unnecessary).

This document itself was drafted by a further Claude (Fable~5) session
from the archived transcripts, with every transcribed statement
re-verified against the sources. Three small items added at write-up
time --- the automatic-quotient bound in Remark~\ref{rem:autow}, the
sharpness construction in Remark~\ref{rem:wall-sharp}, and the $m=3$
example in Remark~\ref{rem:m3} --- were subsequently put through a
targeted cross-family critique of their own (codex; transcript:
\path{provenance/critique-note-additions-codex.md}).
The sharpness construction and the $m=3$ example survived complete
verification, and that run contributed the explicit extremizers now in
Remark~\ref{rem:wall-sharp} and the part-(iii) witness now in
Remark~\ref{rem:m3}. The quotient-constant remark was refuted as
originally drafted --- it presented $\lfloor m/2\rfloor$ as the exact
constant, whereas at $m=4$ the true maximum under $|u|\le4$ is $1$, not
$2$ --- and Remark~\ref{rem:autow} has been corrected per that critique
to state only the (valid) inequalities. Their proofs are one-liners (the $m=3$
example a finite check) and are included in full above. Human
involvement consisted of operating the harness,
selecting the problem, and reviewing this document. The proofs are short
and elementary, and the reader is invited to check them directly.

\begin{thebibliography}{99}

\bibitem{BP22}
J.~Bhanja and R.~K. Pandey,
\emph{On the minimum size of subset and subsequence sums in integers},
C. R. Math. Acad. Sci. Paris \textbf{360} (2022), 1099--1111.
\href{https://doi.org/10.5802/crmath.361}{doi:10.5802/crmath.361}.

\bibitem{EP875}
T.~F. Bloom, \emph{Erd\H{o}s problems}, Problem 875,
\url{https://www.erdosproblems.com/875} (accessed July 2026).

\bibitem{CGKL23}
S.~Cambie, J.~Gao, Y.~Kim and H.~Liu,
\emph{The Erd\H{o}s distinct subset sums problem in a modular setting},
arXiv:2308.03748 (2023).

\bibitem{DF95}
J.-M. Deshouillers and G.~A. Freiman,
\emph{On an additive problem of Erd\H{o}s and Straus, 1},
Israel J. Math. \textbf{92} (1995), 33--43.

\bibitem{DF99}
J.-M. Deshouillers and G.~A. Freiman,
\emph{On an additive problem of Erd\H{o}s and Straus, 2},
Ast\'erisque \textbf{258} (1999), 141--148.

\bibitem{DGMS07}
M.~DeVos, L.~Goddyn, B.~Mohar and R.~\v{S}\'amal,
\emph{A quadratic lower bound for subset sums},
Acta Arith. \textbf{129} (2007), no.~2, 187--195.

\bibitem{ENS91}
P.~Erd\H{o}s, J.-L. Nicolas and A.~S\'ark\H{o}zy,
\emph{Sommes de sous-ensembles},
J. Th\'eor. Nombres Bordeaux \textbf{3} (1991), no.~1, 55--72.

\bibitem{Mazur26}
GPT-5.5 Pro and L.~Mazur,
\emph{On an absolute-gap variant of Erd\H{o}s Problem \#875},
note and Lean formalization, May 2026,
\url{https://github.com/lechmazur/erdos_875}.

\bibitem{MR15}
F.~Monopoli and I.~Z. Ruzsa,
\emph{Carries and the arithmetic progression structure of sets},
arXiv:1506.08869 (2015).

\bibitem{Nathanson95}
M.~B. Nathanson,
\emph{Inverse theorems for subset sums},
Trans. Amer. Math. Soc. \textbf{347} (1995), no.~4, 1409--1418.

\bibitem{Olson68}
J.~E. Olson,
\emph{An addition theorem modulo $p$},
J. Combin. Theory \textbf{5} (1968), 45--52.

\bibitem{Straus66}
E.~G. Straus,
\emph{On a problem in combinatorial number theory},
J. Math. Sci. \textbf{1} (1966), 77--80.

\end{thebibliography}

\end{document}
